
\chapter{总结}
\label{chap:conclusion}
随着Higgs粒子的发现，精确测量该粒子的性质成为粒子物理学家的首要目标之一。为了达成此目标，精确计算Higgs
粒子的各个产生道如胶子聚合道$gg\to H$、VBF道$pp\to Hjj$、矢量玻色子伴随产生道$pp\to VH$和顶夸克对伴随
产生道$pp\to t\bar{t}H$是十分重要的。另一方面，Higgs粒子是电弱自发对称破缺机制的成因，所以研究Higgs势和
Higgs自耦合也是十分必要的。为了重构Higgs势，必须研究Higgs对产生的过程。
此外，为了精确检验电弱自发对称破缺机制，研究规范玻色子的耦合也是十分重要的。
最后，为了探索其它可能的Higgs粒子，研究新物理模型也是十分有意义且有趣的。当然，Higgs粒子性质的精确测量也
离不开对背景过程的精确研究，所以研究这些背景过程将有助于精确测量Higgs粒子的性质。
本文主要研究LHC上$pp \to W^+W^-+{\rm jet} \to e^+ \mu^- \nu_e \bar \nu_{\mu}+{\rm jet}+X$
产生过程，以及2HDM类型II中VBF道产生Higgs粒子对的过程。

本文首先介绍标准模型的粒子谱及构造标准模型和新物理模型所需的定域规范不变性，接着
讨论自发对称破缺与标准模型的Higgs机制。为了研究新物理中的Higgs现象，我们选定双Higgs
二重态模型，主要讨论该模型的Higgs场部分和当前的实验对模型参数的限制。在第二章最后的理论部分我们讨论了
一般性的电弱修正理论，例如电弱修正的阶数问题、电弱参数输入方案的选取、光子诱导的过程以及光子夸克事例的
精确判别等。这些问题是计算电弱修正时经常遇到的。
关于电弱修正重要性的讨论也可以查阅Les Houches TeV国际高能物理会议\cite{LHwishlist}，该会议特别强调了
矢量玻色子对伴随jet产生过程的次领头阶电弱修正的重要性。

第三章讨论了LHC上$pp\to W^+W^-+{\rm jet}+X$过程的次领头阶QCD和电弱修正，
并且考虑了带有自旋关联效应的W玻色子的衰变。该计算是目前为止对该过程最为精确的理论预言。
我们还讨论了夸克混合效应，并利用CKM矩阵的幺正性显著简化了计算。
另外，计算还涉及b夸克的处理和红外安全的jet算法，
我们采用了五味道方案和b夸克标记方法来去除顶夸克共振的贡献。
对于QCD和电弱虚修正的计算，会遇到小Gram行列式带来的数值不稳定问题，
我们开发了双精度和四精度联合计算的程序包来解决数值不稳定问题，且该程序包极大地提高了计算效率。
对电弱次领头阶的实光子辐射计算会涉及到四体末态$WW+jet+\gamma$，
对于jet和光子共线的事例，我们采用光子占jet和光子体系的能量比
$z_{\gamma}$来区分光子事例和jet事例。但为了保证红外安全性，
我们创新性地引入夸克光子碎裂函数来吸收掉额外的共线发散。
最后我们还考虑了W玻色子衰变的自旋关联效应，这样的计算比传统的窄宽度近似要精确很多。
我们采用MadSpin方法处理矢量玻色子自旋关联的衰变，从领头阶的计算可以看出，
自旋关联的效应在特定区域可以达到$\pm 10\%-20\%$。从末态粒子的
微分分布中我们可以得出，当领头jet的横动量截断为$p_{T,j} > 1~TeV$时，电弱相对修正在重整化和因子化能标
取成W玻色子的质量时可以达到-22.2\%，充分证明电弱Sudakov效应在高能量区域特别是TeV量级会变得非常显著且重要。

我们在第四章中讨论了2HDM模型下VBF道产生轻的CP为偶的Higgs对的次次领头阶QCD修正。首先，Higgs粒子对的产生
过程可以用来研究Higgs粒子的自耦合，并且可以重建Higgs势。另外2HDM模型有更加丰富的Higgs物理现象学，此过程
同样可以用于探索新的Higgs粒子。为了取得更高精度的结果，我们采用结构函数法计算次次领头阶QCD修正。
次次领头阶修正显著降低了积分总截面对能标的依赖，表明这一阶计算十分重要。我们还
讨论了$PDF+\alpha_s$组合不确定度，结果表明其与能标不确定度大致相当。我们还展示了总截面对模型
参数的依赖以及$h^0$粒子的微分分布和相空间依赖的K因子，结果表明K因子对应的修正为正且不大于10\%。
当$H^0$粒子质量足够大时，$h^0$粒子对可以通过$H^0$共振产生，所以此过程为探测新的Higgs粒子提供一种可能
的方式。


